% ============================================================
% §2.10 · Uso de Inteligencia Artificial Generativa
% ============================================================

\section{Uso de Inteligencia Artificial Generativa}

Durante el desarrollo del proyecto K-QGMIP se utilizó
\textbf{Claude} (Anthropic) mediante la interfaz \textbf{Claude Code},
a lo largo de las cuatro fases del proyecto (18 mayo -- 15 junio 2026).

El registro completo de todas las interacciones --- prompts, decisiones
técnicas, archivos modificados y resultados de pruebas --- está
disponible en la siguiente carpeta de Google Drive:

\begin{center}
  \href{https://drive.google.com/drive/u/1/folders/1_O1T8vP5PjtyhC1141kj9PLuoT6PGLo8}%
  {\texttt{traceability\_data/} \textemdash{} Registro completo de sesiones IA}
\end{center}

La carpeta contiene 22 archivos Markdown, uno por sesión, que documentan:

\begin{itemize}
  \item \textbf{Etapa y fecha} de cada sesión de trabajo.
  \item \textbf{Prompt exacto} escrito por el estudiante.
  \item \textbf{Acciones ejecutadas} por Claude (archivos creados o modificados).
  \item \textbf{Resultados} de tests y benchmarks obtenidos en esa sesión.
\end{itemize}

\subsection*{Áreas de asistencia}

\begin{table}[H]
\centering
\begin{tabular}{ll}
  \toprule
  \textbf{Área} & \textbf{Uso de IA} \\
  \midrule
  Comprensión del código base & Mapeo de contratos e interfaces heredadas \\
  Validación y corrección & Diseño de suite de tests, detección de bugs \\
  Implementación KQNodes & Clase \texttt{KQNodes}, criterio C4, tests \\
  Implementación KGeoMIP & Clase \texttt{KGeoMIP}, estrategia E4, optimizaciones \\
  Documentación & Manual Técnico y Manual de Usuario \\
  \bottomrule
\end{tabular}
\caption{Áreas del proyecto con asistencia de IA.}
\label{tab:areas_ia}
\end{table}

\subsection*{Reflexión}

La IA actuó como acelerador técnico, pero cada decisión fue revisada
y aprobada por el equipo.
Los errores detectados por Claude (bugs DT-10, DT-11, $\Delta\Phi$
incorrecto) requirieron comprensión matemática del estudiante para ser
validados.
La limitación principal fue la pérdida de contexto entre sesiones,
mitigada mediante la carpeta de trazabilidad referenciada arriba.
